% !TEX root = ../../main.tex

\section{主要研究内容}

围绕基于 HarmonyOS 的听障人士交互训练系统设计与实现，本文主要完成了以下工作。

第一，完成 HearBridge 系统的总体架构设计。
系统采用多端协同架构，由 HarmonyOS 移动端、Spring Boot 服务端、Vue 后台管理端和 Python 手势识别服务组成。
移动端负责课程学习、训练交互、相机采集、视频上传、数字人展示和结果反馈；
服务端负责业务接口、用户认证、数据存储、文件管理、识别服务调用和语义修正；
后台管理端负责训练资源、样本数据和模型版本维护；
Python 识别服务负责关键点提取、模型训练、实时识别和句子视频识别。

第二，完成移动端训练应用的设计与实现。
系统基于 HarmonyOS 平台开发，使用 ArkTS 与 ArkUI 构建课程浏览、课程详情、训练流程、数字人动作示范、实时识别结果展示、视频上传和句子识别结果展示等页面。
移动端围绕用户训练路径组织页面流程，使用户能够从课程学习进入具体训练项，并在训练后查看识别结果和反馈信息。

第三，完成服务端业务接口和后台管理能力的设计与实现。
服务端基于 Spring Boot 构建，结合 MyBatis、MySQL、Redis 和 MinIO 完成用户认证、业务数据管理、文件资源存储和识别服务调用。
后台管理端基于 Vue 3、TypeScript 和 Element Plus 实现，提供手势分类、手势资源、样本管理、训练任务和模型版本管理等功能，使系统具备训练资源维护和识别模型迭代能力。

第四，完成基于关键点特征的手势识别服务设计与实现。
Python 识别服务基于 FastAPI 对外提供接口，使用 MediaPipe 提取手部和姿态关键点，使用 TensorFlow 进行模型训练和推理。
服务支持实时单词识别、raw dataset 样本采集、样本扫描、特征转换、模型训练、模型运行时状态查询和模型重载等能力，为移动端训练和后台模型维护提供支撑。

第五，完成有限词表下的句子视频识别与语义修正链路。
系统支持用户上传手语句子视频，由 Python 识别服务完成视频读取、动作窗口分析、词级候选序列生成和中文映射；
Spring Boot 服务端基于识别结果调用大语言模型进行语义修正和句子组织，最终向移动端返回原始词序列、中文展示和修正后的句子结果。
该部分工作验证了将视觉识别模型与语义后处理方法结合用于受限场景辅助交流的可行性。

第六，完成系统功能测试与实验验证。
本文围绕用户登录、课程浏览、手势训练、实时识别、视频上传识别、后台资源管理、样本管理、模型训练、模型发布和文件资源访问等核心功能进行测试，并结合有限词表句子视频识别实验分析系统在识别链路、语义修正和用户展示方面的实际效果。
